\begin{figure*}[t]
\centering
\hfil\hfil%
\begin{subfigure}[t]{.3\textwidth}
\centering
\scalebox{0.7}{%
\begin{tikzpicture}[
  node distance=2.5em, auto,
  lat/.style={draw=black, circle, minimum size=2em},
  det/.style={draw=black, rectangle, minimum size=2em},
  obs/.style={circle, draw=black, fill=black!20, minimum size=2em},
  gen/.style={->, -{Stealth[length=.5em, inset=0pt]}},
  inf/.style={dashed, ->, -{Stealth[length=.5em, inset=0pt]}},
]
% Observations
\node[obs, inner sep=.02em] (o1) {$o_1,r_1$};
\node[obs, right=of o1, inner sep=.02em] (o2) {$o_2,r_2$};
\node[lat, right=of o2, inner sep=.02em] (o3) {$o_3,r_3$};
% States
\node[det, above=of o1] (s1) {$h_1$};
\node[det, above=of o2] (s2) {$h_2$};
\node[det, above=of o3] (s3) {$h_3$};
% Actions
\node[obs, above=of s1] (a1) {$a_1$};
\node[lat, above=of s2] (a2) {$a_2$};
% Generative
\path (s1) edge[gen] node {} (o1);
\path (s2) edge[gen] node {} (o2);
\path (s3) edge[gen] node {} (o3);
\path (s1) edge[gen] node {} (s2);
\path (s2) edge[gen] node {} (s3);
\path (a1) edge[gen] node {} (s2);
\path (a2) edge[gen] node {} (s3);
% Inference
\path (s1) edge[inf, bend right=30] node {} (s2);
\path (a1) edge[inf, bend left=30] node {} (s2);
\path (o2) edge[inf, bend left=30] node {} (s2);
% Spacing
% \node[det, above=of s3] (p1) {};
% \node[lat, above=of a1] (p2) {};
% \path (p1) edge[gen, bend left=35] (o3);
\end{tikzpicture}}
\caption{Deterministic model (RNN)}
\label{fig:rnn}
\end{subfigure}\hfil%
%
\begin{subfigure}[t]{.3\textwidth}
\centering
\scalebox{0.7}{%
\begin{tikzpicture}[
  node distance=2.5em, auto,
  lat/.style={draw=black, circle, minimum size=2em},
  det/.style={draw=black, rectangle, minimum size=2em},
  obs/.style={circle, draw=black, fill=black!20, minimum size=2em},
  gen/.style={->, -{Stealth[length=.5em, inset=0pt]}},
  inf/.style={dashed, ->, -{Stealth[length=.5em, inset=0pt]}},
]
% Observations
\node[obs, inner sep=.02em] (o1) {$o_1,r_1$};
\node[obs, right=of o1, inner sep=.02em] (o2) {$o_2,r_2$};
\node[lat, right=of o2, inner sep=.02em] (o3) {$o_3,r_3$};
% States
\node[lat, above=of o1] (s1) {$s_1$};
\node[lat, above=of o2] (s2) {$s_2$};
\node[lat, above=of o3] (s3) {$s_3$};
% Actions
\node[obs, above=of s1] (a1) {$a_1$};
\node[lat, above=of s2] (a2) {$a_2$};
% Generative
\path (s1) edge[gen] node {} (o1);
\path (s2) edge[gen] node {} (o2);
\path (s3) edge[gen] node {} (o3);
\path (s1) edge[gen] node {} (s2);
\path (s2) edge[gen] node {} (s3);
\path (a1) edge[gen] node {} (s2);
\path (a2) edge[gen] node {} (s3);
% Inference
\path (s1) edge[inf, bend right=30] node {} (s2);
\path (a1) edge[inf, bend left=30] node {} (s2);
\path (o2) edge[inf, bend left=30] node {} (s2);
% Spacing
% \node[det, above=of s3] (p1) {};
% \node[lat, above=of a1] (p2) {};
% \path (p1) edge[gen, bend left=35] (o3);
\end{tikzpicture}}
\caption{Stochastic model (SSM)}
\label{fig:ssm}
\end{subfigure}\hfil%
%
\begin{subfigure}[t]{.3\textwidth}
\centering
\scalebox{0.7}{%
\begin{tikzpicture}[
  node distance=2.5em, auto,
  lat/.style={draw=black, circle, minimum size=2em},
  det/.style={draw=black, rectangle, minimum size=2em},
  obs/.style={circle, draw=black, fill=black!20, minimum size=2em},
  gen/.style={->, -{Stealth[length=.5em, inset=0pt]}},
  inf/.style={dashed, ->, -{Stealth[length=.5em, inset=0pt]}},
]
% Observations
\node[obs, inner sep=.02em] (o1) {$o_1,r_1$};
\node[obs, right=of o1, inner sep=.02em] (o2) {$o_2,r_2$};
\node[lat, right=of o2, inner sep=.02em] (o3) {$o_3,r_3$};
% States
\node[lat, above=of o1] (s1) {$s_1$};
\node[lat, above=of o2] (s2) {$s_2$};
\node[lat, above=of o3] (s3) {$s_3$};
\node[det, above=of s1] (b1) {$h_1$};
\node[det, above=of s2] (b2) {$h_2$};
\node[det, above=of s3] (b3) {$h_3$};
% Actions
\node[obs, above=of b1] (a1) {$a_1$};
\node[lat, above=of b2] (a2) {$a_2$};
% Generative
\path (s1) edge[gen] node {} (o1);
\path (s2) edge[gen] node {} (o2);
\path (s3) edge[gen] node {} (o3);
\path (b1) edge[gen, bend left=35] node {} (o1);
\path (b2) edge[gen, bend left=35] node {} (o2);
\path (b3) edge[gen, bend left=35] node {} (o3);
\path (b1) edge[gen] node {} (s1);
\path (b2) edge[gen] node {} (s2);
\path (b3) edge[gen] node {} (s3);
\path (b1) edge[gen] node {} (b2);
\path (b2) edge[gen] node {} (b3);
\path (s1) edge[gen] node {} (b2);
\path (s2) edge[gen] node {} (b3);
\path (a1) edge[gen] node {} (b2);
\path (a2) edge[gen] node {} (b3);
% Inference
\path (b2) edge[inf, bend right=30] node {} (s2);
\path (o2) edge[inf, bend left=30] node {} (s2);
\end{tikzpicture}}
\caption{Recurrent state-space model (RSSM)}
\label{fig:rssm}
\end{subfigure}%
\hfil\hfil%
\caption{Latent dynamics model designs.
In this example, the model observes the first two time steps and predicts the third.
Circles represent stochastic variables and squares deterministic variables. Solid lines denote the generative process and dashed lines the inference model.
(a)~Transitions in a recurrent neural network are purely deterministic. This prevents the model from capturing multiple futures and makes it easy for the planner to exploit inaccuracies.
(b)~Transitions in a state-space model are purely stochastic. This makes it difficult to remember information over multiple time steps.
(c)~We split the state into stochastic and deterministic parts, allowing the model to robustly learn to predict multiple futures.}
\label{fig:models}
\end{figure*}
